\documentclass[a4,11pt]{aleph-notas}

% -- Paquetes adicionales
\usepackage{enumitem}
\usepackage{aleph-comandos}

% -- Datos
\institucion{Facultad de Ciencias Exactas, Naturales y Ambientales}
\carrera{Catálogo STEM}
\asignatura{Álgebra lineal}
\tema{Ejercicios 01: Matrices}
\autor{Andrés Merino}
\fecha{Periodo 2026-1}

\logouno[0.16\textwidth]{Logos/logoPUCE_04_ac}
\definecolor{colortext}{HTML}{1957d1}
\definecolor{colordef}{HTML}{1957d1}
\fuente{montserrat}

% -- Comandos adicionales
\setlist[enumerate]{label=\roman*.}

\begin{document}

\encabezado

%%%%%%%%%%%%%%%%%%%%%%%%%%%%%%%%%%%%%%
\section{Matrices}
%%%%%%%%%%%%%%%%%%%%%%%%%%%%%%%%%%%%%%

%%%%%%%%%%%%%%%%%%%%%%%%%%%%%%%%%%%%%%%%
\begin{ejer}
    Dada la matriz 
    \[
        A = \begin{pmatrix} 
            1 & 2 & 3 \\ 
            4 & 5 & 6 \\ 
            -1 & 0 & 2
            \end{pmatrix},
    \]
    calcule:
\begin{enumerate}
    \item $A^\intercal + A$
    \item $A - A$
    \item $A + A$
    \item $3A$
\end{enumerate}
\end{ejer}

\begin{proof}[Solución]\hspace{0pt}
    \begin{enumerate}
    \item Se tiene que:
    \[
    A^\intercal + A 
    = \begin{pmatrix} 1 & 4 & -1\\ 2 & 5 & 0 \\  3 & 6 &2\end{pmatrix} +
      \begin{pmatrix} 1 & 2 & 3 \\ 4 & 5 & 6 \\ -1 & 0 &2\end{pmatrix} 
    = \begin{pmatrix} 2 & 6 & 2 \\ 6 & 10& 6 \\  2 & 6 & 4  \end{pmatrix}
    \]
    \item Se tiene que:
    \[
    A - A 
    = \begin{pmatrix} 1 & 2 & 3 \\ 4 & 5 & 6 \\ -1 & 0 & 2\end{pmatrix}
    - \begin{pmatrix} 1 & 2 & 3 \\ 4 & 5 & 6 \\ -1 & 0 & 2\end{pmatrix}
    = \begin{pmatrix} 0 & 0 & 0 \\ 0 & 0 & 0 \\  0 & 0 & 0\end{pmatrix}
    \]
    \item Se tiene que:
    \[
    A + A 
    = \begin{pmatrix} 1 & 2 & 3 \\ 4 & 5 & 6 \\ -1 & 0 & 2\end{pmatrix}
    + \begin{pmatrix} 1 & 2 & 3 \\ 4 & 5 & 6 \\ -1 & 0 & 2\end{pmatrix}
    = \begin{pmatrix} 2 & 4 & 6 \\ 8 & 10 & 12 \\ -2 & 0 & 4\end{pmatrix}
    \]
    \item Se tiene que:
    \[
    3A 
    = 3\begin{pmatrix} 1 & 2 & 3 \\ 4 & 5 & 6 \\ -1 & 0 & 2\end{pmatrix}
    =  \begin{pmatrix} 3 & 6 & 9 \\ 12 & 15 & 18 \\ -3 & 0 & 6\end{pmatrix}
    \]
\end{enumerate}
\end{proof}

%%%%%%%%%%%%%%%%%%%%%%%%%%%%%%%%%%%%%%%%
\begin{ejer}
    Utilizando las matrices
    \[
    A = \begin{pmatrix}1&2&3&-1\\-2&3&1&5\end{pmatrix}, \quad 
    B=\begin{pmatrix} 1 &0 &2 \\ -1& -1& 1\\ 0&2& 0 \\ 1& 1& 0\end{pmatrix}, \quad C=\begin{pmatrix}1&0&0\\2&1&0 \\3&2&1\end{pmatrix}, 
    \]
    \[
    D=\begin{pmatrix}1&2&3\\0&1&0\\-1&0&2\end{pmatrix} \texty E=\begin{pmatrix}1&2&0&0\\0&0&-1&-2\end{pmatrix};
    \] 
    calcule:
    \begin{enumerate}
        \item $AB$;
        \item $BC$;
        \item $B(C+D)$;
        \item $(E+A)B$.
    \end{enumerate}
\end{ejer}

\begin{proof}[Solución]\hspace{0pt}
\begin{enumerate}
    \item $AB = \begin{pmatrix} -2&3&4\\0&4&-1\end{pmatrix}$.
    \item $BC= \begin{pmatrix}7&4&2\\0&1&1\\4&2&0\\3&1&0\end{pmatrix}$.
    \item $B(C+D)=\begin{pmatrix}1&0&2\\-1&-1&1\\0&2&0\\1&1&0\end{pmatrix}\begin{pmatrix} 2&3&3\\2&2&0\\2&2&3 \end{pmatrix}
    =\begin{pmatrix}6&6&9\\-2&-2&0\\4&4&0\\4&4&3\end{pmatrix}$.
    \item $(E+A)B = \begin{pmatrix}2&4&3&-1\\-2&3&0&3\end{pmatrix}
    \begin{pmatrix}1&0&2\\-1&-1&1&\\0&2&0\\1&1&0\end{pmatrix}
    =\begin{pmatrix}-3&1&8\\-2&0&-1\end{pmatrix}$.
\end{enumerate}
\end{proof}

%%%%%%%%%%%%%%%%%%%%%%%%%%%%%%%%%%%%%%%%
\begin{ejer}
    Considere la matriz
    \[ 
        A=\begin{pmatrix}
            1 & 2 & 0 \\
            0 & 1 & 3 \\
            0 & 0 & 1 
        \end{pmatrix} 
    \]
    y sea $B=A-I$; donde $I$ es la identidad. Calcule $B^{n}$ para todo $n>0$.
\end{ejer}

\begin{proof}[Solución]\hspace{0pt}
Para $n=1$,
    \[
        B^{1}=B=A-I= \begin{pmatrix}
            0 & 2 & 0 \\
            0 & 0 & 3 \\
            0 & 0 & 0
        \end{pmatrix}.
    \]
Para $n=2$, 
    \[
        B^{2}=B  B = \begin{pmatrix}
            0 & 0 & 6 \\
            0 & 0 & 0 \\
            0 & 0 & 0 
        \end{pmatrix}.
    \]
Para $n=3$, 
    \[
        B^{3}=B^{2} B= \begin{pmatrix}
            0 & 0 & 0 \\
            0 & 0 & 0 \\
            0 & 0 & 0
        \end{pmatrix} 
        = 0.
    \]
que es la matriz nula.

Para $n\geq 3$, $B^{n}=B^{3} B^{n-3}=0$.
\end{proof}

%%%%%%%%%%%%%%%%%%%%%%%%%%%%%%%%%%%%%%%%
\begin{ejer}
    Considere las matrices
    \[
        A=\begin{pmatrix}1&-2&0\\3&1&4\end{pmatrix},\quad
        B=\begin{pmatrix}2&1\\0&-1\\5&3\end{pmatrix},\quad
        C=\begin{pmatrix}4&0&-1\\2&3&1\end{pmatrix}.
    \]
    Determine cu\'ales de las siguientes operaciones est\'an definidas y calcule su resultado:
    \begin{enumerate}
        \item $A+C$;
        \item $AB$;
        \item $BA$;
        \item $B+C$.
    \end{enumerate}
\end{ejer}

\begin{proof}[Soluci\'on]\hspace{0pt}
\begin{enumerate}
    \item $A+C$ est\'a definida porque ambas matrices son de orden $2\times 3$:
    \[
        A+C=
        \begin{pmatrix}1&-2&0\\3&1&4\end{pmatrix}+
        \begin{pmatrix}4&0&-1\\2&3&1\end{pmatrix}
        =\begin{pmatrix}5&-2&-1\\5&4&5\end{pmatrix}.
    \]
    \item $AB$ est\'a definida porque $A$ es $2\times 3$ y $B$ es $3\times 2$:
    \[
        AB=
        \begin{pmatrix}1&-2&0\\3&1&4\end{pmatrix}
        \begin{pmatrix}2&1\\0&-1\\5&3\end{pmatrix}
        =\begin{pmatrix}2&3\\26&14\end{pmatrix}.
    \]
    \item $BA$ tambi\'en est\'a definida porque $B$ es $3\times 2$ y $A$ es $2\times 3$:
    \[
        BA=
        \begin{pmatrix}2&1\\0&-1\\5&3\end{pmatrix}
        \begin{pmatrix}1&-2&0\\3&1&4\end{pmatrix}
        =\begin{pmatrix}5&-3&4\\-3&-1&-4\\14&-7&12\end{pmatrix}.
    \]
    \item $B+C$ no est\'a definida porque $B$ es de orden $3\times 2$ y $C$ es de orden $2\times 3$.
\end{enumerate}
\end{proof}

%%%%%%%%%%%%%%%%%%%%%%%%%%%%%%%%%%%%%%%%
\begin{ejer}
    Halle $x,y,z\in\R$ para que las matrices sean iguales:
    \[
        \begin{pmatrix}
            2x-1 & y+3 & 4\\
            z-2 & 5 & x+y
        \end{pmatrix}
        =
        \begin{pmatrix}
            5 & 1 & 4\\
            0 & 5 & 9
        \end{pmatrix}.
    \]
\end{ejer}

\begin{proof}[Soluci\'on]\hspace{0pt}
Por igualdad de matrices, las entradas correspondientes deben ser iguales:
\[
2x-1=5,\qquad y+3=1,\qquad z-2=0,\qquad x+y=9.
\]
De las tres primeras ecuaciones se obtiene
\[
x=3,\qquad y=-2,\qquad z=2.
\]
Sin embargo,
\[
x+y=3+(-2)=1\neq 9.
\]
Por lo tanto, el sistema es incompatible y \textbf{no existen} $x,y,z\in\R$ que satisfagan la igualdad.
\end{proof}

%%%%%%%%%%%%%%%%%%%%%%%%%%%%%%%%%%%%%%%%
\begin{ejer}
    Sean
    \[
        A=\begin{pmatrix}1&2\\-1&0\\3&4\end{pmatrix},\quad
        B=\begin{pmatrix}0&5\\2&-3\\1&1\end{pmatrix},\quad
        \alpha=-2.
    \]
    Calcule:
    \begin{enumerate}
        \item $(A+B)^\intercal$ y $A^\intercal + B^\intercal$;
        \item $(\alpha A)^\intercal$ y $\alpha A^\intercal$.
    \end{enumerate}
\end{ejer}

\begin{proof}[Soluci\'on]\hspace{0pt}
\begin{enumerate}
    \item Primero,
    \[
        A+B=
        \begin{pmatrix}1&7\\1&-3\\4&5\end{pmatrix}
        \implies
        (A+B)^\intercal=
        \begin{pmatrix}1&1&4\\7&-3&5\end{pmatrix}.
    \]
    Adem\'as,
    \[
        A^\intercal=\begin{pmatrix}1&-1&3\\2&0&4\end{pmatrix},
        \qquad
        B^\intercal=\begin{pmatrix}0&2&1\\5&-3&1\end{pmatrix},
    \]
    y por tanto
    \[
        A^\intercal+B^\intercal=
        \begin{pmatrix}1&1&4\\7&-3&5\end{pmatrix}.
    \]
    As\'i, $(A+B)^\intercal=A^\intercal+B^\intercal$.

    \item
    \[
        \alpha A=-2A=
        \begin{pmatrix}-2&-4\\2&0\\-6&-8\end{pmatrix}
        \implies
        (\alpha A)^\intercal=
        \begin{pmatrix}-2&2&-6\\-4&0&-8\end{pmatrix}.
    \]
    Por otro lado,
    \[
        \alpha A^\intercal=-2\begin{pmatrix}1&-1&3\\2&0&4\end{pmatrix}
        =\begin{pmatrix}-2&2&-6\\-4&0&-8\end{pmatrix}.
    \]
    Luego, $(\alpha A)^\intercal=\alpha A^\intercal$.
\end{enumerate}
\end{proof}

\end{document}
